\PARstart{A}{rtificial} intelligence systems are now deployed in settings where errors carry
consequences for safety, livelihoods, and fundamental rights: hiring, credit, clinical
decision support, critical infrastructure, and increasingly autonomous, tool-using
agents. In response, a converging body of governance instruments now asks developers and
deployers of such \emph{high-risk} systems to \emph{demonstrate}, rather than merely
assert, that their systems are trustworthy. In the United States, the National Institute
of Standards and Technology (NIST) AI Risk Management Framework (AI RMF~1.0) establishes a
voluntary but widely referenced structure for identifying, measuring, and managing AI
risk~\cite{nist2023airmf}, complemented by a Generative AI Profile~\cite{nist2024genai}.
In the European Union, the AI Act, Regulation (EU)~2024/1689, imposes binding
obligations on high-risk systems, including risk management, data governance, technical
documentation, logging, transparency, human oversight, and accuracy, robustness, and
cybersecurity~\cite{euaiact2024}. Internationally, ISO/IEC~42001:2023 specifies a
certifiable AI management system with controls spanning the AI life
cycle~\cite{iso42001}. Together these frameworks define the conformity surface that a
high-risk AI system must satisfy.

A recurring word in all three instruments is \emph{evidence}: conformity is established
by demonstrable, documented results, not by good intentions. Yet the frameworks are
deliberately written at the level of \emph{outcomes} and \emph{processes} (``the system
shall achieve appropriate levels of accuracy and robustness,'' ``fairness and bias shall
be evaluated and managed'') and do not prescribe \emph{which} technical evaluation
methods produce the required evidence. In parallel, a large and fast-moving technical
literature has produced concrete methods for evaluating AI systems: fairness and bias
benchmarks, adversarial-robustness tests, red-teaming and jailbreak suites,
interpretability and explanation-faithfulness analyses, calibration and uncertainty
quantification, data and model documentation, human-oversight studies, and post-deployment
monitoring. These two worlds, the regulatory and the technical, have largely developed
in isolation. A practitioner asking the operational question, ``\emph{which evaluation do
I run to evidence this obligation?},'' finds no systematic answer in the engineering
literature.

\textbf{Contribution.} This article provides that answer. Prior work has mapped regulatory
concepts to \emph{standards} and compared frameworks pairwise (Section~\ref{sec:related});
our contribution is, to our knowledge, the first \emph{unified} crosswalk of concrete AI
evaluation methods to the \emph{simultaneous} requirements of all three frameworks (NIST
AI RMF, the EU AI Act, and ISO/IEC~42001) in the technical and engineering literature retrieved by our search, together with a
practitioner playbook for assembling cross-framework conformity evidence. Our central
question is operational: which technical evaluation evidence supports conformity with which
NIST AI RMF function and subcategory, which EU AI Act high-risk obligation (Articles
9--15), and which ISO/IEC~42001 Annex~A control. The signature artifact is a
\emph{methods $\times$ requirements mapping matrix} (Table~\ref{tab:matrix},
Fig.~\ref{fig:heatmap}) in which each cell records the technical relevance (primary,
supporting, or contextual) that a method family supplies for a requirement, grounded in the
official clause text and the verified primary literature. The mapping is a structured,
auditable argument graded by technical relevance, not a deterministic compliance verdict. From the matrix we derive a
\emph{gap analysis} (Section~\ref{sec:gaps}) that identifies requirements weakly served
or unserved by current technical methods, including a structural \emph{measurement gap}
for agentic systems whose behavior, rather than whose token-level outputs, must be
assessed.

\textbf{Positioning.} Existing treatments of AI regulation are predominantly legal or
comparative: jurisdictional overviews of global governance
requirements~\cite{global2024governance}, framework-to-framework crosswalks, and
standards-concept mappings~\cite{hernandez2024knowledgegraph}. Valuable as these are,
none bridges from \emph{technical evaluation methods} to regulatory requirements at the
level an engineer can act on. Conversely, technical surveys of evaluation
methods~\cite{chang2024survey,yehudai2025survey} catalogue methods without reference to
conformity obligations. We build the engineering bridge between the two, and we
state that position explicitly. We also situate two prior, unpublished
frameworks (one for structured high-stakes evaluation and one for per-component agent
fairness, citations withheld for double-blind review) as instances within the broader
method taxonomy, mentioned as context rather than as independent evidence and therefore
not listed in the references.

\textbf{Scope and framing.} We focus on high-risk and high-stakes AI, including large
language models (LLMs) and agentic systems, and we adopt a US-forward framing that leads
with the NIST AI RMF and the US assurance ecosystem before turning to the EU and
international instruments. The taxonomy comprises twelve method families; eleven are
established, while privacy evaluation (M12) is included for completeness but is comparatively
emerging and thinly covered, as we note in its analysis. We follow a structured, transparent review protocol for a reproducible,
auditable synthesis, and we verify cited works and regulatory claims against official
sources where publicly accessible. This article is a technical review and does \emph{not} constitute legal
advice; the mapping is an engineering interpretation intended to guide evaluation
practice.

\textbf{Roadmap.} Section~\ref{sec:background} introduces the three frameworks and their
shared lifecycle spine. Section~\ref{sec:related} positions the work.
Section~\ref{sec:method} details the review methodology. Section~\ref{sec:taxonomy}
presents the taxonomy of twelve technical evaluation method families.
Section~\ref{sec:mapping} presents the mapping, the central contribution.
Section~\ref{sec:gaps} analyzes the gaps. Section~\ref{sec:discussion} translates the
mapping into a practitioner conformity-evidence playbook. Sections~\ref{sec:limits}
through~\ref{sec:conclusion} discuss limitations, open challenges, and conclusions.
